\documentclass[class=NCU_thesis, crop=false]{standalone}
\begin{document}

\chapter{Abstract}
According to statistics released by Taiwan’s Ministry of Health and Welfare in early 2026, the number of registered individuals with hearing impairment in Taiwan has exceeded 140,000, and Taiwan Sign Language (TSL) serves as an important medium for their daily communication. However, most existing online sign language resources are organized as isolated word lookup systems or fixed sentence videos. When users wish to express complete Chinese sentences, they often need to search for corresponding sign videos word by word and combine them manually, resulting in fragmented and discontinuous information. To improve the usability of digital sign language content, this study develops a modular system that converts Chinese text input into virtual human sign language animation by integrating Chinese-to-gloss conversion, sign motion data retrieval, motion parameter concatenation, and avatar rendering.

The system design separates sign motion data construction from user-requested animation generation. During sign motion data construction, EHM-Tracker converts existing sign language videos into parameterized tracking data containing body pose, hand motion, and facial information, and a gloss-based index is established for motion retrieval. During animation generation, a front-end Chinese-to-natural-sign-language translation module, namely the Sign Language Translation (SLT) module, generates a gloss sequence from Chinese input, after which the system retrieves and validates the tracked motion data corresponding to each gloss. Motion envelope analysis is then applied to extract the effective motion interval of each sign word. Optional temporal smoothing is applied to word-level motion parameters, and transition frames are inserted between adjacent sign segments to reduce pose discontinuities and motion jumps during multi-word concatenation. Finally, the concatenated motion parameters are passed into the GUAVA avatar rendering module to generate virtual human sign language animation with consistent character appearance.

In the experiments, word-level smoothing was evaluated on 10 samples. The normalized jerk cost of SMPL-X body pose and MANO hand pose decreased by an average of 99.90\% and 99.87\%, respectively, indicating a substantial reduction in high-frequency motion variation. However, the average hand-only PCK@0.10 decreased from 0.9476 to 0.8868, revealing a trade-off between motion smoothness and frame-level hand alignment. This study also compares EHM-Tracker results under 100-step and 1000-step optimization settings. The results show that 1000-step optimization generally improves reprojection error and PCK when both hands are clearly visible. However, in single-hand samples with low visibility or hands moving out of frame, longer optimization may increase the risk of mistracking. Therefore, this study adopts the more conservative 100-step setting for constructing the sign motion database. In addition, a user questionnaire was conducted to evaluate 25 system-generated videos, with 53 valid responses collected. The average scores of the four evaluation metrics ranged from 3.80 to 4.01, with accuracy and gesture clarity receiving relatively higher scores. The questionnaire results indicate that the proposed system is feasible for generating virtual human sign language animation from Chinese text, while hand details, facial expressions, and sentence-level naturalness remain areas for future improvement. Overall, this study demonstrates that a modular pipeline combining SLT, existing sign language video resources, EHM-Tracker, and GUAVA can serve as a foundation for future applications in accessible communication, sign language learning, and digital sign language content generation.

\vspace{2em}
\noindent \textbf{Keywords:} \keywordsEn{} % Set keywords in config.tex
\end{document}
